\section{1. Introduction}\label{dependence-risk-in-a-reallocation-with-unchanged-acquisition-cost}

How much of a portfolio's diversification benefit survives uncertainty about dependence? A risk desk may have agreed marginal models for its individual positions without an equally well-determined model for their joint losses with the background book \cite{pflugpohl2018}. We ask how a proposed reallocation performs when these marginal models are retained but the coupling changes. Specifically, we seek the smallest dependence perturbation that reverses its Expected Shortfall (ES) advantage when both allocations face the same stressed law.

Let \(X_1,X_2\) be the unit losses of two positions of equal initial cost. A transfer from holdings \((1+\varepsilon,1-\varepsilon)\) to \((1,1)\), with an unchanged background loss \(\beta X_3\), changes the portfolio loss from \(A_\varepsilon=(1+\varepsilon)X_1+(1-\varepsilon)X_2+\beta X_3=B+\varepsilon H\) to \(B\), where
\[
B=X_1+X_2+\beta X_3,\qquad H=X_1-X_2,\qquad \beta>0.
\]
At upper-tail probability \(p\in(0,1)\), we compare them through \(\Phi_\varepsilon(P)=ES_p^P(B)-ES_p^P(A_\varepsilon)\), optimizing the two ES thresholds separately under \(P\) \cite{rockafellar2002}, so that \(\Phi_\varepsilon(P)<0\) favors the transfer under that law. Although \(\beta X_3\) cancels from the payoff difference \(A_\varepsilon-B\), its dependence with the traded positions determines which states enter each portfolio's upper tail. Put differently, a trade assessed in isolation can rank differently once embedded in the whole book \cite{brandtner2018}.

Writing \(\Pi\) for the Fr\'echet class of joint laws with these marginals, we measure distance from the reference law \(P_0\) by weighted \(W_1\) with ground cost \(c_\gamma(x,y)=|y-x|_\gamma=\sum_i\gamma_i|y_i-x_i|\), where each \(\gamma_i>0\) specifies the budget penalty per unit coordinate loss displacement as losses are re-paired across joint states. At budget \(\rho\), the stress problem is
\[
D_\varepsilon(\rho)
=\sup_{\substack{P\in\Pi\\ W_\gamma(P,P_0)\le\rho}}
\bigl\{ES_p^P(B)-ES_p^P(A_\varepsilon)\bigr\},
\qquad m_\varepsilon=-D_\varepsilon(0).
\]
Under an exchangeable reference law, the traded positions have equal marginal ES contributions, so the nominal gain is of second order: \(m_\varepsilon=(J/p)\varepsilon^2+o(\varepsilon^2)\), with \(J>0\). For example, when \(\gamma_1\le\min\{\gamma_2,\gamma_3/\beta\}\), the local exchange in Proposition~\ref{prop:nominal-benchmark} attains an initial stress derivative \(K_\varepsilon\) of order one, suggesting the reversal budget \(r_{0,\varepsilon}=m_\varepsilon/K_\varepsilon=O(\varepsilon^2)\). The derivative at zero budget does not determine whether this efficiency persists until the nominal gain is exhausted. Since the predicted reversal already occurs at scale \(\rho=s\varepsilon^2\), we need the risk profile over a nonvanishing interval of the rescaled budget \(s\), rather than its tangent at zero alone.

Three effects contribute at the critical scale. Moving the \(O(\varepsilon)\) mass in an \(O(\varepsilon)\)-layer around the reference tail boundary by distances of order \(\varepsilon\) changes both transport cost and tail expectation by \(O(\varepsilon^2)\), matching the nominal gain. This local movement need not preserve the marginals. For a fixed regular local plan, the marginal CDF discrepancies are \(O(\varepsilon^2)\); by Vallender's formula \cite{vallender1974}, their \(L^1\) norms give the one-dimensional restoration costs, even when compensation occurs far from the tail boundary. An \(O(\varepsilon)\) shift of either ES threshold also contributes at order \(\varepsilon^2\) through the curvature of its objective. We must therefore optimize the local displacement, marginal restoration, and both thresholds together.

When the reference law has a positive continuous density on a compact box \(I^2\times I_3\) and is exchangeable in its first two coordinates, Theorem~\ref{thm:critical-limit} gives a finite, continuous profile
\[
F(s)=\lim_{\varepsilon\downarrow0}\varepsilon^{-2}D_\varepsilon(s\varepsilon^2),
\qquad F(0)=-J/p,
\]
with convergence locally uniform in \(s\ge0\). The transport energies have a \(\Gamma\)-limit \cite{dalmaso1993} that includes both the local displacement cost and the cost of restoring its marginal projections. In the recovery argument, coordinatewise quantile maps attain the latter cost for each fixed regular local plan. Although this correction moves losses throughout the box, its leading effects on the two ES values cancel outside the threshold layer. Approximation then extends the paired-risk convergence to kernels with finite displacement moment.

Let \(\rho^-_\varepsilon\) and \(\rho^+_\varepsilon\) be the infima of budgets permitting a tie and a strict reversal. The strict-growth theorem identifies their common leading coefficient as the unique zero \(s^*\) of \(F\):
\[
\rho^-_\varepsilon=s^*\varepsilon^2+o(\varepsilon^2),\qquad
\rho^+_\varepsilon=s^*\varepsilon^2+o(\varepsilon^2).
\]
\emph{Initial efficiency and its limit.} The effective unit costs are \(\gamma_1,\gamma_2\), and \(c_3=\gamma_3/\beta\), with \(c_{\min}=\min\{\gamma_1,\gamma_2,c_3\}\). If the first traded coordinate is cheapest, Proposition~\ref{prop:balanced-linear-profile} constructs a positive interval \([0,s_{\rm lin}]\) on which \(F(s)=-J/p+s/(p\gamma_1)\). To reach zero at this rate, however, efficient displacement must remain possible throughout the required interval. Theorem~\ref{thm:zero-premium} characterizes attainment of the universal efficiency bound:
\[
s^*=Jc_{\min}\quad\Longleftrightarrow\quad
c_3\le\min\{\gamma_1,\gamma_2\}.
\]
We call the nonnegative gap \(s^*-Jc_{\min}\) the \emph{marginal-restoration premium} (normalized by \(Jc_{\min}\) when reported as a percentage). When background displacement is cheapest, exchangeability permits opposite shifts to balance on every fixed-background slice. When only traded directions are cheapest, a corner region requires a displacement moment of order \(\delta^2\), but its coordinate projections admit only \(O(\delta^4)\) of compensating moment near the diagonal (Figure~\ref{fig:corner}). Under the attained-slope condition of Proposition~\ref{prop:nominal-benchmark}, with \(c_{\min}=\gamma_1<c_3\), this gives \(\rho^\pm_\varepsilon/r_{0,\varepsilon}\to s^*/(\gamma_1J)>1\).

\emph{An analytical benchmark.} Consider independent uniform losses on \([0,2]^3\), with \(\beta=1\), \(p=1/2\), and \(\gamma=(1,1,2)\). Here \(g=3/8\) and \(J=5/32\), so Proposition~\ref{prop:nominal-benchmark} gives \(m_\varepsilon=(5/16)\varepsilon^2+o(\varepsilon^2)\). By Corollary~\ref{prop:uniform-variable-cost}, the profile is exactly \(F(s)=2s-5/16\) for \(0\le s\le1/96\), yet its zero satisfies \(13/64\le s^*\le29/96\). Even in this uniform model, the proved lower bound on reversal is already \(30\%\) above the local-efficiency prediction \(5/32\), despite an attainable initial slope.

\emph{Sensitivity at equal effective costs.} Along the cost family \(\gamma_1=\gamma_2=\gamma\), \(c_3=\gamma(1+\eta)\) for \(\eta\ge0\), the baseline \(Jc_{\min}=\gamma J\) is constant, so \(s^*\) and the marginal-restoration premium have the same right derivative \(\gamma Ja_+\) at \(\eta=0\). Here \(a_+\in(0,1)\) is the least background share among efficient stresses, characterized in Theorem~\ref{thm:projected-flow} by continuous maximum-flow/minimum-cut duality; Section 5 shows how this share and the budget sensitivity \(\gamma Ja_+\) can move in opposite directions as the reference dependence changes.

Our question lies between extremal risk aggregation and local distributional sensitivity. Subadditivity bounds aggregate risk \cite{artzner1999,mcneil2015}, while rearrangements within a Fr\'echet class describe how diversification depends on the coupling \cite{ruschendorf2013,embrechts2013,bernard2014,embrechts2015,wang2015superadditive}. Partial information leads to convex-order and mixability constraints \cite{li2018partial}, but can still permit comonotonic upper tails \cite{devecchinendelstreicher2026}, and fixed-marginal MES bounds quantify uncertainty in individual tail contributions \cite{chenfurmanlin2026}. For allocation comparisons, the symmetry imposed on the candidate laws is decisive. Chen et al.\ \cite{chen2022} and Liu and Liu \cite{liu2026} use exchangeability and majorization to identify diversification protection. Liu and Liu also analyze when complete dependence ambiguity favors concentration and how Wasserstein uncertainty modifies that conclusion. In our model, stressed couplings need not remain exchangeable. Symmetry at the reference gives stationarity, which does not by itself protect the reallocation once symmetry is broken under stress. We use this setting to measure the distance to a specified decision reversal, in the spirit of reverse sensitivity analysis \cite{pesenti2019,pesenti2022}.

Transport duality provides a natural way to preserve the marginal information \cite{villani2003,santambrogio2015}. In Wasserstein robust optimization \cite{blanchet2019,gao2023}, adding marginal constraints confines the perturbation to dependence changes \cite{gao2017,eckstein2020}. Such constraints encode calibrated risk-neutral laws in multiasset robust pricing \cite{eckstein2021multiasset}; Wasserstein and copula restrictions also arise in benchmark portfolio optimization \cite{pesentijaimungal2023}. Here the constraint acts on a common stressed loss law for two allocations. After rescaling, its full coordinate projections remain as a transport penalty, and the efficient allocation of displacement reduces to a capacity-constrained flow problem \cite{korman2015}.

Local Wasserstein expansions relate the sensitivity of optimized decisions and AVaR criteria to changes in the reference law \cite{bartl2021}, with adapted transport incorporating the information available in multiperiod decisions \cite{bartlwiesel2023}. When a first-order term vanishes, higher-order expansions become relevant. Fiechtner and Blanchet \cite{fiechtner2026} derive a second-order regret expansion under a unique empirical-risk minimizer in this situation, and Sauldubois \cite{sauldubois2026} develops higher-order Wasserstein and adapted-Wasserstein expansions, including martingale constraints. Dependence restrictions also affect local sensitivity, as in Lam \cite{lam2018}, while Namkoong \cite{namkoong2026} evaluates the first-order response over restricted transport directions. In our problem the trade size and stress radius vanish together, so locating a decision reversal at scale \(\rho=s\varepsilon^2\) requires the full rescaled profile, with global marginal compensation and both optimized ES thresholds retained over the budget interval needed to reach zero.

Section 2 establishes the nominal gain and the initial stress slope. Section 3 derives the variational limit at scale \(\rho=s\varepsilon^2\) and identifies the unique zero of the limiting profile. Section 4 characterizes zero premium and represents the boundary derivative by continuous projected flow--cut duality. Section 5 develops analytical examples and their implications for reverse stress testing. Appendix A contains the analytic proofs, and Appendix B gives the potential estimates and section calculations for the examples.

\paragraph{Notation and conventions} We write \(T_\#\mu\) for a pushforward and measure displacements by \(|d|_\gamma=\sum_i\gamma_i|d_i|\). The constant in \(O_\alpha(\varepsilon^k)\) may depend on the fixed kernel and model, but not on \(\varepsilon\). Scalar section integrands are extended by zero outside their domains, and projection identities are equalities of finite signed Borel measures. We use \(\inf\varnothing=+\infty\); allocation ratios are assigned the value zero where both numerator and denominator vanish.
